\section{The Next Instrumental Sub-goal Is Encoded Before Execution}
\label{sec:what}

Our first question is whether the next instrumental goal $g_k$ is actually present and readable at the \pivot $\pi_k$, \emph{before} any of it is written. A model that improvises token by token would only reveal its goal as it decodes; a model that decides up front would already carry a representation of \textit{``I am about to do $g_k$''} at $\pi_k$. We test this by \emph{reading out} the model's goal at the pivot and checking whether $g_k$ appears in the read-out.

\subsection{Approach}
\label{sec:what:approach}

\paragraph{\goalscope.} 
% We define \emph{intent} as the \emph{future step} the LM implements in response to a prompt at the current token position. By \emph{step} here we mean a coherent concept (\emph{``checking if a query is valid''}), not the specific next token (or sequence of tokens). In our controlled setup ({\color{red} Figure \Cref{fig:teaser}}), the sub-task at (or in near future) of the current token position is the target \emph{intent}.\footnote{\small For natural code-generation, beyond our controlled setup, the target \emph{intent} is not always cleanly separable as it is not clear which hierarchy level we should care about. A detailed discussion and goalscope read-out results are available on \Cref{app:unspecified}.}

We find that we can read-out LM's intent from the latents at a single token position, with a simple verbalization approach. Our verbalization approach, \goalscope is a lightweight tool that is essentially an extension of \emph{Patchscope} presented at \cite{ghandeharioun2024patchscope}, with 2 important distinctions: (i) instead of residual stream latents, we locate a set of attention heads and use what they contribute to the residual stream; (ii) we check multi-token rollouts, with interventions at every generated tokens, instead of checking a single token. See {\color{red} Figure \Cref{fig:goalscope-setup}}\arnab{todo} for a schema of the setup. We leave the details in Appendix \ref{app:goalscope}.

\paragraph{Locating \goal heads.} 
\arnab{a lot of text here. make more concise. details should be discussed in the Appendix.}
We locate the \goal heads with activation patching (\Cref{sec:notations:interventions}) on a simpler prompt setup. Both the source and destination prompt asks the LM to implement a single predicate ($\mathrm{pred}_\mathrm{src} \neq \mathrm{pred}_\mathrm{dst}$) on the same data type (int | str | $\cdots$). Both the prompts are cut-off at a token position, such that the next generated token is a representative of the predicate being implemented. Assume that $\mathrm{tok}_\mathrm{src}$ and $\mathrm{tok}_\text{dest}$ are the next tokens the LM predicts with the source and destination prompts.

We aim to locate the representations in the source run $\mathrm{M}(p_\mathrm{src})$ that encode the abstract concept/predicate being implemented in code. To do that, we \emph{patch} the representations, cached from the source run, to the model run with the destination prompt, and track the logit/prob of the token $\mathrm{tok}_\mathrm{targ}$. The token $\mathrm{tok}_\mathrm{targ}$ is a representative of if the LM were implementing the predicate $\mathrm{pred}_\mathrm{src}$, instead of $\mathrm{pred}_\mathrm{dest}$, as the next token with the destination prompt. In order to ensure that the localized latents capture the concept, not shallower token-level information, we make certain that $\mathrm{tok}_\mathrm{targ} \neq \mathrm{tok}_\mathrm{src}$ by using different programming languages for source and destination prompts. A concrete example of the contrastive setup is illustrated on {\color{red} Figure \Cref{fig:goal_head_loc}}.

We localize the concept-carrying attention heads by learning a sparse binary mask, jointly over all the attention heads. From the source run, we cache what individual attention heads \emph{write} to the residual stream (their OV contribution,  notation from \cite{elhage2021mathematical}, $a_t$ in \Cref{eq:head}). Then, in the patched run, with the destination prompt, we perform the following intervention per head. 

\vspace{-1.5em}
\begin{align}
    \label{eqn:opt-int}
    a^{\ell j}_{-1} & \leftarrow \mathrm{mask}^{\ell j} * a^{\ell j}_{\mathrm{src}} + (1 - \mathrm{mask}^{\ell j}) * a^{\ell j}_{\mathrm{-1}}
\end{align}
\vspace{-1.5em}

Here, $a^{\ell j}_{i}$ denotes what the head $[\ell, j]$ contribute to the residual stream at $i^\text{th}$ token position (-1 for last token). The terms $\mathrm{mask}^{\ell j})$ are binary values per head, jointly learned to maximize the $\mathrm{logit}(\mathrm{tok}_\mathrm{targ})$ in the patched run.

\subsection{Results}
We apply the \goalscope at different token positions of the text \emph{generated by the model}. To keep the discussion scoped for this section, here we only analyze read-out results at the \pivot tokens. See {\color{red} Figure \Cref{fig:goal-eval}} \arnab{pivot setup with \goalscope faithfulness results}. We have two measures \emph{faithfulness} and \emph{causality}. Success in both of these measures contribute towards our claim that \emph{a coherent intent is formed well before the actual implementation}.

\paragraph{Faithfulness.} We ask if the readout is \emph{faithfully} describing a non-trivial portion of what the LM has actually implemented in the next \emph{meaningful} step (the sub-task $g_k$). In our controlled corpus, we prefill until $g_k$, which naturally reduces the next sub-task to $g_k$, giving us a gold target to measure the fidelity of our readouts.

\paragraph{Causality.} \arnab{the read-out itself is kind of causal. maybe we don't need to check for this again.} We transplant the pivot goal encoding into an information-less code stub in a different programming language and test whether it drives the stub to implement a non-trivial part of $g_k$. Higher causality metric shows that the verbalized goal is the encoding the LM acts on, it is not an encoding tighly couples with the next token distribution.

We take the intent representation in patch it (with the setup at {\color{red} Figure \Cref{fig:goalscope-setup}}) to an information-less code prefix from a different programming language (original implementation in python, read-out probe code-stub in C++). Then we check if the intent encoding can drive the information-less code-stub to implement a non-trivial part of the target $g_k$.

Both \emph{faithfulness} and \emph{causality} requires us to evaluate the concept captured in the multi-token read-out (intent verbalization or code completion) against a gold description (instruction span in the source context). We rely on a Oracle LM (\texttt{claude-sonnet-4.6}) for this purpose (details in \Cref{app:goalscope}).

{\color{red} Figure \Cref{fig:goal-eval}} shows that the read-outs achieves high faithfulness and causality well before the LM generates the first token of the target $g_k$. We include qualitative examples of read-outs on free-form code-generation in \cref{app:unspecified}. 


% We test this with a causal mediation analysis setup. From a source prompt $p_\text{src}$ with a list of sub-tasks $G_\text{src}$, we cache the LM latents at the pivot token position for the $i^\text{th}$ sub-task $g_i \in G_\text{src}$. Then, on the LM run for a different prompt $p_\text{dst}$ with subtasks $G_\text{dst}$ (such that $G_\text{src} \cap G_\text{src} = \emptyset$), we intervene with the cached latents at the pivot token position (and then onwards) for the $j^{th}$ ($i \neq j$) subtask $g_j \in G_\text{dst}$. Check {\color{red} Figure \ref{fig:goal-cause}} for a concrete example. We only patch the outputs for heads used in the \goalscope setup.

% On the patched run, we stop the generation after a single line of code has been generated, and then check if what that code functionally implements $g_i$, the sub-task from the source prompt for which the latents were cached. We use a larger LLM as judge to automate the process for scaling up (see Appendix \ref{app:task} for details). 

\paragraph{Goals heads show \emph{gaze} behavior.} \cite{gandikota2026gaze} observe that while generating the caption of an image, specific attention heads in the LM-backbone gaze (looks back to) the specific segments of the image currently being described. We observe that \goal heads also show gaze behavior: at the pivot token (query token is $\pi_k$) the collective value-weighted attention pattern of \goal heads concentrate on specific tokens in the prompt describing the next sub-task ($g_k$). In \Cref{app:gaze} we quantify whether this gaze behavior aligns with $g_k$ and find high correlational evidence. This gaze behavior, together with the readout, becomes part of our debugging repertoire utilized in later sections to get a better understanding on what our interventions are actually modifying in the latent space, without generating a single token.

\paragraph{Goals heads retain the intent across multiple token positions.}
\dbau{How does the intent representation change across multiple tokens? Check cosine affinity with respect to a \pivot point. Qualitative shifts in the read-out}
\arnab{Not in the main text. Move to a separate sub-section in the appendix. This could be a good segway before the boundary encoding, if we decide to include it.}

\subsection{Discussion}
\label{sec:what:discussion}

The readout results establish that a coherent, verbalizable encoding of the next instrumental goal is presnt at the pivot, before the LM commits to concrete implementation tokens. That the encoding can be readout in natural language indicates that LLMs plan their next step in language. The remaining sections focus on \emph{where} the encoding lives, and \emph{how} it comes to be present at the pivot. \arnab{we don't have an answer to \emph{how} yet.}